%%%%%%%%%%%%%%%%%%%%%%%%%%%%%%%%%%%%%%%%%%%%%%%%%%%%%%%%%%%%%%%%%%%%%%
%
%  T-RO submission draft (IEEEtran, journal mode)
%
%  All experiment-dependent numbers are macros in the PLACEHOLDER block
%  below and render as red \ph{...} until the hybrid-controller runs
%  are finalized. Figures are placeholder boxes with final captions.
%
%  Compile:  pdflatex main_tro && bibtex main_tro && pdflatex x2
%
%%%%%%%%%%%%%%%%%%%%%%%%%%%%%%%%%%%%%%%%%%%%%%%%%%%%%%%%%%%%%%%%%%%%%%
\documentclass[journal]{IEEEtran}

\usepackage{amsmath,amssymb}
\usepackage{graphicx}
\usepackage{booktabs}
\usepackage{multirow}
\usepackage{xcolor}
\usepackage[hidelinks]{hyperref}
\usepackage{url}

\graphicspath{{figures/}}

% ---------------- placeholder machinery ----------------
\newcommand{\ph}[1]{\textcolor{red}{#1}}          % pending number
\newcommand{\phfig}[2]{\fbox{\parbox[c][#1][c]{0.96\linewidth}{\centering
  \textcolor{red}{[Placeholder]}\\ \footnotesize #2}}}
% ---- numbers to be filled from the hybrid-controller runs ----
\newcommand{\nTrain}{18{,}432}
\newcommand{\nDev}{2{,}048}
\newcommand{\nSealed}{10{,}000}
\newcommand{\phOracleGap}{\ph{$+$XX}}        % enumerated vs generative oracle [mm]
\newcommand{\phHeadroomG}{\ph{XX}}           % generative-pool headroom [mm]
\newcommand{\phHeadroomE}{\ph{XXX}}          % enumerated-pool headroom [mm]
\newcommand{\phSpreadG}{\ph{XXX}}            % generative within-task spread [mm]
\newcommand{\phSpreadE}{\ph{XXX}}            % enumerated within-task spread [mm]
\newcommand{\phTransferCap}{\ph{$-$XX}}      % transferred/geometric capture [%]
\newcommand{\phCapSingle}{\ph{XX.X}}         % selector capture, single [%]
\newcommand{\phCapEns}{\ph{XX.X}}            % selector capture, ensemble [%]
\newcommand{\phDeltaVal}{\ph{$+$XX.X}}       % main delta, validation [mm]
\newcommand{\phDeltaExt}{\ph{$+$XX.X}}       % main delta, external [mm]
\newcommand{\phDeltaSealed}{\ph{$+$XX.X}}    % main delta, sealed [mm]
\newcommand{\phCIVal}{\ph{$[\cdot,\cdot]$}}
\newcommand{\phCIExt}{\ph{$[\cdot,\cdot]$}}
\newcommand{\phCISealed}{\ph{$[\cdot,\cdot]$}}
\newcommand{\phTrimVal}{\ph{$+$XX.X}}
\newcommand{\phTrimExt}{\ph{$+$XX.X}}
\newcommand{\phTrimSealed}{\ph{$+$XX.X}}
\newcommand{\phSeedStd}{\ph{X.X}}            % across-training-seed std [mm]

\begin{document}

\title{Extreme Motion Generation via Enumerated Seed Selection\\
and Hybrid Null-Space Control}

\author{Xinyi~Yuan,
        Weiwei~Wan,~\IEEEmembership{Member,~IEEE,}
        and~Kensuke~Harada,~\IEEEmembership{Fellow,~IEEE}%
\thanks{The authors are with the Graduate School of Engineering Science,
The University of Osaka, Japan.}}

\markboth{IEEE Transactions on Robotics}%
{Yuan \MakeLowercase{\textit{et al.}}: Extreme Motion Generation via
Enumerated Seed Selection and Hybrid Null-Space Control}

\maketitle

\begin{abstract}
Industrial operations such as surface coating and welding require a
fixed-base manipulator to follow a prescribed path for as long as possible
before a safety constraint---a joint limit or a tool-orientation
tolerance---forces termination. We refer to the maximization of this
constrained path length as extreme motion generation and address it with a
two-stage system. Before execution, initial joint configurations are
enumerated on the cone-constrained inverse-kinematics solution set,
diversified across the disconnected branches of the self-motion manifold,
and a learned selector commits to a single configuration; during execution,
a hybrid controller resolves the redundancy, delegating the well-sampled
interior of the configuration space to a reinforcement-learning policy and
the near-boundary region to a classical null-space controller through a
hysteresis rule. The selector is supervised by the complete closed-loop
return of every feasible candidate, and we show that this supervision is
necessary: on the enumerated candidate set, selection rules built from
instantaneous geometric criteria perform worse than a trivial
first-feasible rule, whereas the return-supervised selector recovers
approximately two thirds of the attainable headroom. We further locate the
limit of any execution-free selection scheme, showing that an exact
kinematic continuation of each candidate adds no selective information
beyond the learned features. On \nSealed{} straight-line path-following
tasks with a 7-DoF manipulator, the proposed system lengthens the average
executed path by \phDeltaSealed~mm over a learned generative initialization
baseline under an identical execution budget of one selector inference and
one rollout per task, and the improvement is consistent across independent
training seeds and a sealed final evaluation.
\end{abstract}

\begin{IEEEkeywords}
Extreme motion generation, redundant manipulators, seed selection,
reinforcement learning, null-space control.
\end{IEEEkeywords}

\IEEEpeerreviewmaketitle

%======================================================================
\section{Introduction}
\label{sec:intro}
\IEEEPARstart{C}{onsider} a fixed-base manipulator coating or welding a
large workpiece along a straight line. Intuition suggests that the
attainable stroke is bounded by the reachable workspace; in practice,
however, the motion terminates much earlier, when a joint reaches its
limit or the tool axis leaves its orientation tolerance. The practically
achievable path length is therefore governed by how long these safety
boundaries can be deferred during execution. We call the problem of
maximizing the constrained Cartesian path length along a prescribed
trajectory \emph{extreme motion generation}. Its industrial value is
direct: the length of a single continuous motion determines the area a
manipulator can process in one setup, and exploiting the kinematic
capability more fully substitutes for enlarging the robot or relocating
its base.

Two decisions determine the outcome of such a motion, and they act on
different time scales. The first is made before the motion begins: for a
task that constrains fewer degrees of freedom than the arm possesses, the
feasible initial configurations form a self-motion manifold that joint
limits partition into disconnected branches, and configurations on
different branches realize the same end-effector pose while supporting
very different downstream path lengths. The second decision is made at
every control step: the redundancy must be resolved so that the commanded
joint motion defers the safety boundaries rather than drifting toward
them. This paper presents a system that addresses both decisions and, in
doing so, characterizes how much each one can contribute.

The initialization stage raises a question that, to our knowledge, has
not been answered quantitatively: \emph{how much of a long-horizon motion
outcome can be predicted from information available before the motion
begins?} The question is nontrivial for three reasons. First, the choice
is decisive: on our task distribution the best and the worst feasible
starting configuration of the same task differ in final path length by a
median of \phSpreadE~mm, comparable to the average task length itself.
Second, the consequence of the choice is revealed late: whether a branch
survives is dominated by orientation-cone violations that occur deep into
the trajectory and correlate only weakly with any instantaneous geometric
quantity evaluated at the start. Third, the decision must be made without
trial. A budget of one execution per task excludes probing, and because a
complete episode spans only a few tens of control steps, any probe long
enough to be informative would cost nearly as much as executing every
candidate exhaustively. Initialization is therefore forced to compress
the outcome of an entire closed-loop execution into a single static
decision, and the compression loss can be measured.

The execution stage poses a complementary trade-off between foresight and
robustness. A classical null-space controller optimizes instantaneous
secondary objectives; it is myopic but predictable, and remains reliable
arbitrarily close to the joint-limit boundary. A learned policy possesses
the long-horizon foresight that the objective demands, but its training
data concentrate in the interior of the configuration space, and its
behavior degrades exactly in the near-boundary region that extreme motion
generation must traverse. We reconcile the two by partitioning the
configuration space by proximity to the joint-limit boundary and
assigning each region to the controller that is reliable there, with a
hysteresis rule suppressing chattering at the interface.

Concretely, the proposed system operates as follows. For each task,
cone-constrained inverse kinematics is solved from randomized restarts,
producing tens of distinct solutions that are deduplicated and reduced by
farthest-point selection to a diverse candidate set spanning the
self-motion branches. A permutation-equivariant selector network, trained
on the complete closed-loop return of every feasible candidate under the
frozen controller, scores the set once and commits to a single
configuration. The hybrid controller then executes the motion to
termination. The deployed budget is one selector inference, one seed, and
one rollout per task; no candidate is probed and no model-based rollout
is performed at execution time.

The contributions are as follows. First, we show that the candidate set
itself, rather than the selection rule, can bound such a system: direct
enumeration of the cone-constrained IK solution set raises the attainable
per-task ceiling by \phOracleGap~mm on average over a learned generative
proposal distribution and roughly doubles the selectable headroom, with
the ceiling saturating once the enumeration covers the distinct
self-motion branches. Second, we establish that supervision by complete
closed-loop returns is necessary for selection on such a set: every
instantaneous geometric criterion we evaluate, and every selector trained
on the narrower generative proposals, performs worse than taking the
first feasible candidate, whereas the identical feature set supervised by
complete returns recovers about two thirds of the headroom. Third, we
measure where execution-free selection ends: enlarging the model,
extending its training, and even granting the selector an exact kinematic
continuation of every candidate all leave the recovered fraction
unchanged, indicating that the residual headroom depends on the learned
policy's closed-loop behavior and is invisible to any static computation.
Fourth, we verify that, under the proposed training distribution and
task-relative observation design, the execution stage requires no
adaptation to the selection-induced initial-state distribution, and we
provide a stratified analysis explaining why. All findings are validated
on \nDev-task development sets and a preregistered, sealed \nSealed-task
final evaluation.

%======================================================================
\section{Related Work}
\label{sec:related}

\subsection{Manipulability, Reachability, and Initialization}
Classical redundancy resolution optimizes instantaneous dexterity
measures such as the manipulability
ellipsoid~\cite{yoshikawa1985manipulability}, and modern planners embed
these measures into sampling strategies and cost
functions~\cite{jin2017manipulability,shen2023adaptive,cai2025just}.
Reachability analysis complements local dexterity with global feasibility
envelopes~\cite{yang2009placement,makhal2018reuleaux}, which have been
extended toward predictive and safety-oriented
variants~\cite{feng2025predictive,jauhri2022robot,selim2022safe}. These
criteria act on the instantaneous state; the present work quantifies
their limitation on a long-horizon objective, showing that on a diverse
candidate set they are not merely suboptimal but counterproductive as
selection rules.

\subsection{Learned Initialization and Generative IK}
Learning-based methods approximate the multi-modal solution set of
inverse kinematics with normalizing flows~\cite{ames2022ikflow} or
diffusion models~\cite{zhang2025ikdiffuser}, and seed downstream planners
or controllers with sampled
configurations~\cite{huang2024diffusionseeder,yoon2023learning}. Such
generators concentrate probability mass on regions that were favorable in
their training data, which is precisely the property we find limiting:
the proposals cluster within a narrow quality band, so the attainable
ceiling of any system built on them is capped. We compare against a
conditional-diffusion initialization baseline of this type and show that
replacing generation with explicit enumeration of the constraint set
raises the ceiling substantially while removing a trained component.

\subsection{Combining Learned and Model-Based Control}
Residual formulations superpose a learned correction on a classical
controller~\cite{johannink2019residual,silver2018residual}, while
switching formulations select among controllers according to safety
criteria~\cite{mehmood2022black,alshiekh2018safe,kurtz2021control}. Our
execution stage follows the switching paradigm, but the trigger is not
imminent constraint violation alone: the configuration space is
partitioned by proximity to the joint-limit boundary so that each
controller operates where it is statistically reliable, with the
objective of advancing the motion rather than merely remaining safe.

\subsection{Sequential Policy Chaining}
When a task decomposes into stages, the output distribution of one stage
becomes the input distribution of the next, and bidirectional fine-tuning
has been proposed to align adjacent policies~\cite{chen2023sequential}.
We adopt the same credit structure across our two
stages---initialization must be evaluated by the success of
execution---but treat the necessity of each adaptation direction as an
experimental question. In our system the backward direction (returns
supervising selection) carries the entire measured benefit, while the
forward direction (controller adaptation to the selection-induced
distribution) is verified to be unnecessary, a result we attribute to the
coverage induced by task-relative observations and randomized training
tasks.

%======================================================================
\section{Problem Statement}
\label{sec:problem}
Let $\mathbf{q}\in\mathbb{R}^{n}$ denote the joint configuration of an
$n$-DoF manipulator, $\mathbf{p}(\mathbf{q})\in\mathbb{R}^{3}$ the
end-effector position, and $\mathbf{z}(\mathbf{q})\in\mathbb{R}^{3}$ the
tool axis. A straight-line path-following task is specified by the
condition vector
\begin{equation}
\mathbf{c}=[\mathbf{p}_{0},\mathbf{d},\mathbf{n}]\in\mathbb{R}^{9},
\end{equation}
where $\mathbf{p}_{0}$ is the start position and $\mathbf{d}$,
$\mathbf{n}$ are unit vectors specifying the motion direction and the
nominal tool axis. The end-effector must track a reference that advances
along the line at speed $v$ while satisfying, at every step $t$ until
termination,
\begin{align}
\|\mathbf{p}(\mathbf{q}_{t})-\mathbf{p}_{t}^{\mathrm{ref}}\|
  &\leq \epsilon_{\mathrm{pos}},\label{eq:c-pos}\\
\arccos\!\big(\mathbf{z}(\mathbf{q}_{t})^{\!\top}\mathbf{n}\big)
  &\leq \theta_{\max},\label{eq:c-cone}\\
\mathbf{q}_{\min}\preceq \mathbf{q}_{t} &\preceq
  \mathbf{q}_{\max}.\label{eq:c-jl}
\end{align}
The episode terminates at the first violation, and the objective is the
net path length $\mathcal{L}$ accumulated along $\mathbf{d}$ before
termination.

Because only the three-dimensional position is tracked exactly, the
manipulator is kinematically redundant for $n>3$, and two couplings
determine $\mathcal{L}$. The initial configuration $\mathbf{q}_{0}$ must
be chosen from the cone-constrained IK solution set of
$(\mathbf{p}_{0},\mathbf{n})$, a set whose joint-limit-induced branches
are disconnected in configuration space; and the redundancy must be
resolved at every subsequent step. We formalize the two decisions as a
two-stage optimization
\begin{equation}
\max_{\phi,\theta}\;\;
\mathbb{E}_{\mathbf{c}}\Big[
\mathcal{L}\big(\mathbf{q}_{i},\,\pi_{\theta};\,\mathbf{c}\big)\Big],
\qquad
i\sim\pi^{\mathrm{sel}}_{\phi}\big(\cdot\mid\mathbf{c},
Q(\mathbf{c})\big),
\label{eq:objective}
\end{equation}
where $Q(\mathbf{c})=\{\mathbf{q}_{1},\dots,\mathbf{q}_{K}\}$ is a finite
candidate set produced from the constraint manifold,
$\pi^{\mathrm{sel}}_{\phi}$ is the selection policy, and $\pi_{\theta}$
is the execution controller. Deployment is restricted to a single static
pass: $\pi^{\mathrm{sel}}_{\phi}$ is evaluated once per task, commits to
one $\mathbf{q}_{0}$, and $\pi_{\theta}$ executes one episode. No
candidate probing or model-predictive rollout is permitted at execution
time. This budget reflects the application setting and also isolates the
scientific question of Sec.~\ref{sec:intro}: any information the
selection uses must be computable before motion begins.

%======================================================================
\section{System Overview}
\label{sec:overview}
Fig.~\ref{fig:system} summarizes the system. Offline, a training pool of
tasks is processed by the candidate layer
(Sec.~\ref{sec:enumeration}), every feasible candidate is executed to
termination under the frozen controller, and the resulting complete
return tables supervise the selector (Sec.~\ref{sec:selector}). Online,
the candidate layer and one selector inference reduce each task to a
single initial configuration, from which the hybrid controller
(Sec.~\ref{sec:controller}) executes the motion. The interfaces are
deliberately narrow: the selector communicates with the controller only
through the chosen $\mathbf{q}_{0}$, and the controller communicates with
the selector only through the return labels gathered offline.

\begin{figure}[!t]
\centering
\phfig{4.2cm}{System overview. Top: deployed path---task
$\mathbf{c}$ $\rightarrow$ cone-IK enumeration $\rightarrow$ candidate
set $Q(\mathbf{c})$ ($K{=}32$ + fallback) $\rightarrow$
return-supervised set selector $\rightarrow$ single $\mathbf{q}_{0}$
$\rightarrow$ hybrid controller $\rightarrow$ path length
$\mathcal{L}$; badge: one selector inference, one seed, one rollout, no
probes. Bottom: offline training---complete rollouts of all feasible
candidates under the frozen hybrid controller produce return tables that
supervise the selector (backward credit); the interior policy is trained
by PPO on randomized tasks with mixed reset distributions.}
\caption{Overview of the proposed two-stage system for extreme motion
generation.}
\label{fig:system}
\end{figure}

%======================================================================
\section{Seed Candidate Enumeration}
\label{sec:enumeration}
The candidate layer approximates the cone-constrained IK solution set of
each task by direct enumeration. Tool directions are sampled uniformly on
the spherical cap of half-angle $\theta_{\max}$ about $\mathbf{n}$, and
for each direction a damped least-squares IK problem is solved from mixed
random restarts (uniform and joint-limit-biased), yielding $128$ attempts
per task. Solutions are validated against joint limits with margin,
position error, the orientation cone, and collision, then deduplicated at
$0.08$~rad in joint space; the median task retains on the order of $50$
distinct solutions distributed over the disconnected self-motion
branches. Farthest-point selection in joint space reduces the set to
$K{=}32$ candidates, and a deterministic classical fallback configuration
is appended as a designated final slot, giving the selector a defined
abstention target. Every retained slot passes strict physical
validation before entering the action set.

Two properties of this construction matter downstream. Diversity is the
operative quantity: under farthest-point ordering the attainable ceiling
saturates once the branches are covered
(Sec.~\ref{sec:res-pool}), so additional candidates beyond that point add
redundancy rather than headroom. And enumeration is exhaustive rather
than learned: on the small fraction of tasks where the $K{=}32$ set lacks
a competitive candidate, densifying the enumeration closes most of the
remaining gap, indicating that no generative prior is required at any
point in the pipeline.\footnote{Throughout, ``feasible'' means passing all validation
checks of this section; ``first feasible'' denotes the first valid slot
in enumeration order and serves as the selection-free reference.}

%======================================================================
\section{Return-Supervised Seed Selection}
\label{sec:selector}

\subsection{Supervision Signal}
With the execution controller frozen, every feasible candidate of every
training task is rolled out to termination, producing a complete
per-candidate return table $\{\mathcal{L}_{jk}\}$ over tasks $j$ and
candidates $k$. Complete tables, rather than sampled interactions, make
the supervision exhaustive over the discrete action set: the residual
error of any selector trained on them measures what its input features
cannot express, a property used in Sec.~\ref{sec:res-boundary}.

\subsection{Features and Architecture}
Each candidate is described by a 45-dimensional feature vector comprising
the controller's initial observation at that configuration, the lateral
ray error, log positional manipulability, and ten task-aligned
differential-kinematics descriptors (the damped line-tracking joint
velocity, its norm, a linearized joint-limit travel horizon, and
directional manipulability). The selector is an ensemble of five
permutation-equivariant members. Each member embeds the $K$ candidates,
mean-pools a set context, concatenates it back to every candidate, and
outputs two heads: a listwise score trained by softmax cross-entropy on
within-task range-normalized returns, and a feasibility regression head
trained on returns in metres under a Huber loss. Members are trained on
task-level bootstrap resamples and aggregated by mean score; deployment
is the masked $\arg\max$ over feasible candidates. Feature normalization
uses feasible-slot moments of the training split only.

\subsection{Training Protocol}
The selector is trained on \nTrain{} tasks with complete return tables.
All architectural and schedule decisions are made on a geometry-disjoint
internal split; the development and sealed evaluation sets of
Sec.~\ref{sec:setup} are never read during development. Training a full
ensemble requires minutes on one GPU; the dominant offline cost is the
one-time construction of the return tables.

%======================================================================
\section{Hybrid Null-Space Controller}
\label{sec:controller}

\subsection{Task Decomposition}
The commanded joint velocity separates the tracking task from the
redundancy:
\begin{equation}
\dot{\mathbf{q}}=\mathbf{J}^{+}\dot{\mathbf{x}}
+\big(\mathbf{I}-\mathbf{J}^{+}\mathbf{J}\big)\,\boldsymbol{\xi},
\label{eq:nullspace}
\end{equation}
where $\mathbf{J}$ is the position Jacobian, $\mathbf{J}^{+}$ its damped
pseudo-inverse, $\dot{\mathbf{x}}=v\,\mathbf{d}
+\mathbf{K}_{P}(\mathbf{p}^{\mathrm{ref}}-\mathbf{p})$ advances the
end-effector along the line while rejecting lateral drift, and
$\boldsymbol{\xi}$ is the secondary command projected onto the null
space. Both execution regimes below share this decomposition and differ
only in how $\boldsymbol{\xi}$ is generated, which renders them
interchangeable at any control step.

\subsection{Classical Near-Boundary Regime}
The classical regime performs gradient ascent on a weighted combination
of redundancy objectives,
\begin{equation}
\boldsymbol{\xi}=k_{\mu}\nabla_{\mathbf{q}}w_{\mathbf{d}}(\mathbf{q})
-k_{\mathrm{jl}}\,
\frac{\mathbf{q}-\mathbf{q}_{\mathrm{mid}}}
     {\mathbf{q}_{\max}-\mathbf{q}_{\min}}
+k_{\theta}\,g(\theta)\,
\nabla_{\mathbf{q}}\big(\mathbf{z}(\mathbf{q})^{\!\top}\mathbf{n}\big),
\label{eq:secondary}
\end{equation}
combining directional manipulability $w_{\mathbf{d}}$, joint centering,
and a cone term gated by $g(\theta)$ near the orientation tolerance. Its
behavior is instantaneous and therefore myopic, but it is well defined
and stable arbitrarily close to the joint-limit boundary.

\subsection{Learned Interior Regime}
The interior regime is a PPO-trained policy whose observation exposes
only the quantities required to resolve the redundancy: the normalized
configuration and its element-wise square, the task vectors, the tool
axis and its cone alignment, the lateral ray error, and the previous
action. The end-effector position is deliberately excluded; lateral
tracking is handled by the task term of \eqref{eq:nullspace}, and the
task-relative observation makes states of different tasks commensurable,
a property that Sec.~\ref{sec:res-coupling} shows to be consequential.
The action parameterizes null-space coordinates in a task-aligned
orthonormal basis, so that policy outputs and classical commands live in
the same space. The reward is purely the clipped per-step progress along
$\mathbf{d}$; no shaping toward the boundaries is applied. Training
draws tasks from a pool of $10^{5}$ random lines and mixes reset
configurations among selector samples, uniform feasible candidates, and
the classical fallback.

\subsection{Hysteresis Switching}
At every step the active regime is chosen by the normalized joint-limit
proximity $\eta(\mathbf{q})=\max_{i}
|q_{i}-q_{\mathrm{mid},i}|/q_{\mathrm{half},i}$ with a hysteresis pair
$\tau_{\mathrm{exit}}<\tau_{\mathrm{enter}}$: control passes to the
classical regime when $\eta$ exceeds $\tau_{\mathrm{enter}}$ and returns
to the learned regime only after $\eta$ falls below
$\tau_{\mathrm{exit}}$. The hysteresis suppresses chattering at the
interface, and because both regimes emit commands through
\eqref{eq:nullspace}, switching is transient-free with respect to the
tracking task.

%======================================================================
\section{Experimental Setup}
\label{sec:setup}
All experiments use straight-line path-following with a simulated 7-DoF
Franka FR3 and the constraint constants of
Sec.~\ref{sec:problem}. Tasks are drawn from a pool of $10^{5}$ random
lines filtered for feasibility. The training set comprises \nTrain{}
tasks with complete candidate return tables; two development sets of
\nDev{} tasks each (a validation split, and an external set built from a
freshly seeded pool with audited zero geometry overlap) support the
reported analyses; and a sealed set of \nSealed{} tasks was generated
from fresh seeds after all design decisions, weights, and a
selector-publication rule (the median run across training seeds) were
frozen and preregistered. The sealed set is read once, by the published
selector only.

\subsubsection*{Baselines}
The primary baseline is a learned generative initialization: a
conditional diffusion model samples eight proposals per task, which are
projected onto the exact pose constraint, ranked by a learned selector of
identical capacity and feature set trained on that pool, and executed by
the same hybrid controller. Additional selection references on the
enumerated pool include the first-feasible rule, single geometric
criteria, and the complete-candidate oracle, which executes every
feasible candidate and is reported as a diagnostic upper reference rather
than a method.

\subsubsection*{Metrics}
Systems are compared by paired per-task differences on identical tasks
under the identical controller: the mean difference with a
bootstrap 95\% confidence interval over $2\times10^{4}$ resamples, the
5\% two-sided trimmed mean, and the fractions of tasks changed by more
than $\pm1$~mm. Selection quality within a pool is summarized by the
captured fraction of the oracle headroom,
$(\mathcal{L}_{\mathrm{sel}}-\mathcal{L}_{\mathrm{first}})/
 (\mathcal{L}_{\mathrm{oracle}}-\mathcal{L}_{\mathrm{first}})$.

%======================================================================
\section{Results}
\label{sec:results}

\subsection{The Candidate Set Bounds the System}
\label{sec:res-pool}
Table~\ref{tab:pool} compares the enumerated candidate set against the
generative proposal set on the training tasks under the identical frozen
controller. Enumeration raises the mean attainable ceiling by
\phOracleGap~mm and widens the selectable headroom from \phHeadroomG~mm
to \phHeadroomE~mm per task; the within-task spread between the best and
worst feasible candidate widens from a median of \phSpreadG~mm to
\phSpreadE~mm. Fig.~\ref{fig:pool} traces the ceiling against the pool
size under farthest-point ordering: the curve passes the full
generative-pool ceiling within the first few candidates and saturates
once the self-motion branches are covered, after which additional
candidates contribute at the level of measurement noise. The mechanism
is that generative proposals concentrate in a narrow quality band around
previously favorable regions, whereas enumeration exposes the outcome
diversity of the branches themselves.

\begin{table}[!t]
\caption{Candidate Sets Under the Identical Frozen Controller}
\label{tab:pool}
\centering
\begin{tabular}{lcc}
\toprule
 & generative ($K{=}8{+}1$) & enumerated ($K{=}32{+}1$)\\
\midrule
feasible candidates / task & \ph{--} & \ph{--}\\
first-feasible [m] & \ph{--} & \ph{--}\\
oracle [m] & \ph{--} & \ph{--}\\
headroom [mm] & \phHeadroomG & \phHeadroomE\\
within-task spread, median [mm] & \phSpreadG & \phSpreadE\\
\bottomrule
\end{tabular}
\end{table}

\begin{figure}[!t]
\centering
\phfig{3.4cm}{Curve: mean best path length vs.\ pool size $K$
(farthest-point prefix of the enumerated set), with horizontal reference
lines for the generative-pool oracle and first-feasible level.}
\caption{Attainable ceiling versus candidate-set size.}
\label{fig:pool}
\end{figure}

\subsection{Complete-Return Supervision is Necessary}
\label{sec:res-supervision}
Table~\ref{tab:supervision} isolates the supervision signal while holding
the candidate set, features, and controller fixed. Selection rules built
from instantaneous information---single geometric criteria, and the
generative-pool selector transferred unchanged---capture a negative
fraction of the headroom, i.e., they underperform the first-feasible
rule. The identical feature set supervised by complete closed-loop
returns captures \phCapSingle\% (single model) and \phCapEns\%
(ensemble), with sub-percent variation across training seeds; even a
linear model on the same features and labels captures a substantial
fraction. Fig.~\ref{fig:scaling} reports the capture as a function of
the training-set size. We conclude that on a candidate set diverse
enough to matter, the selection problem is not a ranking refinement but
is created and solved by the supervision signal itself.

\begin{table}[!t]
\caption{Selection Rules on the Identical Enumerated Set}
\label{tab:supervision}
\centering
\begin{tabular}{llc}
\toprule
selection rule & supervision & capture [\%]\\
\midrule
first feasible & none & $0$\\
best geometric criterion & none & \phTransferCap\\
transferred generative-pool selector & generative returns & \phTransferCap\\
linear model, 45-D features & complete returns & \ph{$+$XX}\\
listwise selector (single) & complete returns & \phCapSingle\\
listwise selector (ensemble) & complete returns & \phCapEns\\
\bottomrule
\end{tabular}
\end{table}

\begin{figure}[!t]
\centering
\phfig{3.4cm}{Held-out capture [\%] vs.\ number of training tasks
(log axis), single-model mean $\pm$ std over three seeds and the
member-ensemble curve.}
\caption{Capture of the oracle headroom versus training-set size
(geometry-disjoint internal split).}
\label{fig:scaling}
\end{figure}

\subsection{The Limit of Execution-Free Selection}
\label{sec:res-boundary}
The captured fraction plateaus near two thirds, and the plateau is a
property of the information rather than of the model. Enlarging the
network, doubling the training schedule, and exchanging the encoder
leave the held-out capture within noise of the base configuration. The
decisive control grants the selector the strongest execution-free
information available: for every candidate we integrate the damped
tracking field of \eqref{eq:nullspace}--\eqref{eq:secondary} on the pure
kinematic model until a constraint is met, obtaining an exact nonlinear
continuation of each starting branch without any environment step or
policy call. As an individual criterion this kinematic continuation
rivals the entire trained selector in rank correlation with the true
outcome; appended to the feature set and retrained, it changes the
captured fraction by nothing (Table~\ref{tab:boundary}). The learned
selector had evidently internalized the kinematic continuation from its
one-step descriptors, and the residual headroom is attributable to the
learned interior policy's closed-loop behavior---how it spends the null
space in ways no fixed field anticipates---which is observable only by
execution. This bounds every execution-free selection scheme on this
problem, the present one included.

\begin{table}[!t]
\caption{Probing the Static Information Limit}
\label{tab:boundary}
\centering
\begin{tabular}{lcc}
\toprule
configuration & capture, single [\%] & ensemble [\%]\\
\midrule
base features & \phCapSingle & \phCapEns\\
$+$ capacity / schedule variants & \ph{--} & \ph{--}\\
$+$ kinematic continuation features & \ph{--} & \ph{--}\\
\midrule
kinematic continuation alone ($\arg\max$) &
\multicolumn{2}{c}{\ph{--}}\\
\bottomrule
\end{tabular}
\end{table}

\subsection{Does the Controller Require Adaptation?}
\label{sec:res-coupling}
Sequential-chaining practice suggests adapting the execution stage to the
initial-state distribution induced by the selection
stage~\cite{chen2023sequential}. We tested this direction twice under a
preregistered promotion gate: a naive continuation of PPO on the
selection-induced reset distribution, and a conservatively engineered
round with critic-only warm-up, no entropy bonus, a tightened trust
region, and a fourfold step budget. Neither produced a significant
improvement on either development set, the first being significantly
harmful, and the training reward remained saturated throughout,
indicating that under its own objective the interior policy had nothing
to learn from these states. A stratified analysis
(Fig.~\ref{fig:novelty}) explains the null result and renders it
falsifiable: although the selector chooses configurations far from
anything the generative baseline proposes (median joint-space novelty of
several radians), the adaptation effect is null in every novelty
quartile, and the unadapted policy performs best from the most novel
branches. We attribute this to the combination of task-relative
observations and training over randomized tasks: in the state space the
policy actually observes, an unfamiliar branch of one task coincides
with familiar states of others, so the distribution shift that
adaptation would repair has been absorbed by design. The delineation
also predicts where adaptation should become necessary---small task
pools, absolute observations, or physical hardware---which we leave to
future work.

\begin{figure}[!t]
\centering
\phfig{3.2cm}{Two panels. Left: controller-adaptation effect [mm] by
quartile of joint-space distance between the selected seed and the
nearest generative-pool configuration, with 95\% intervals spanning
zero in all quartiles. Right: path length of the unadapted controller
per novelty quartile, increasing with novelty.}
\caption{Novelty-stratified test of controller adaptation on the
selection-induced distribution.}
\label{fig:novelty}
\end{figure}

\subsection{System Comparison}
\label{sec:res-main}
Table~\ref{tab:main} reports the complete system against the generative
initialization baseline on identical tasks under the identical hybrid
controller and execution budget. The mean executed path lengthens by
\phDeltaVal~mm on the validation set (95\% CI \phCIVal) and
\phDeltaExt~mm on the external set (95\% CI \phCIExt); the trimmed means
of \phTrimVal{} and \phTrimExt~mm indicate that the improvement is
carried by typical tasks rather than by a heavy tail, and the paired
distribution is shown in Fig.~\ref{fig:delta}. Across six independent
selector training seeds the mean gain varies by \phSeedStd~mm. On the
sealed \nSealed-task set, read once with the published selector, the gain
is \phDeltaSealed~mm (95\% CI \phCISealed, trimmed \phTrimSealed~mm).
The deployed cost is unchanged where it matters: one selector inference
(sub-millisecond per task when batched) and one rollout; candidate
enumeration adds a sub-second, execution-free kinematic computation per
task. A fraction of tasks regresses relative to the baseline
(Table~\ref{tab:main}); the enumerated set's weak first-feasible anchor
makes selection errors costlier than on the generative set, and we
discuss the implied risk-sensitive objective in
Sec.~\ref{sec:discussion}.

\begin{table}[!t]
\caption{System Comparison Under Identical Controller and Budget}
\label{tab:main}
\centering
\begin{tabular}{lccc}
\toprule
 & validation & external & sealed\\
 & (\nDev) & (\nDev) & (\nSealed)\\
\midrule
baseline [m] & \ph{--} & \ph{--} & \ph{--}\\
proposed [m] & \ph{--} & \ph{--} & \ph{--}\\
$\Delta$ mean [mm] & \phDeltaVal & \phDeltaExt & \phDeltaSealed\\
95\% CI [mm] & \phCIVal & \phCIExt & \phCISealed\\
$\Delta$ trimmed [mm] & \phTrimVal & \phTrimExt & \phTrimSealed\\
win / regress [\%] & \ph{--} & \ph{--} & \ph{--}\\
capture [\%] & \ph{--} & \ph{--} & \ph{--}\\
\bottomrule
\end{tabular}
\end{table}

\begin{figure}[!t]
\centering
\phfig{3.2cm}{Histogram of per-task paired path-length differences
(proposed $-$ baseline), log-scaled counts, with the mean and the
5\% trimmed mean marked.}
\caption{Distribution of paired per-task differences against the
generative initialization baseline.}
\label{fig:delta}
\end{figure}

%======================================================================
\section{Discussion and Limitations}
\label{sec:discussion}
Three boundaries of the present study merit emphasis. First, the
measured plateau of execution-free selection implies that further gains
on this problem must either price execution-time information below its
oracle cost---an open question, given that informative probes approach
full-episode length---or accept a different deployment budget. Second,
although the mean and trimmed improvements are strongly positive, a
minority of tasks regresses relative to the baseline; the selection
objective optimizes the expected return and is indifferent to per-task
regret, and a risk-sensitive variant that penalizes catastrophic
selections remains unexplored. Third, all findings are obtained in
simulation on a single task family, and the coverage argument of
Sec.~\ref{sec:res-coupling} predicts, but does not verify, that
controller adaptation becomes necessary under transfer to physical
hardware, where the execution distribution shifts for reasons no reset
mixture anticipates.

%======================================================================
\section{Conclusion}
\label{sec:conclusion}
This paper addressed extreme motion generation with a two-stage system:
enumeration of the cone-constrained IK solution set followed by
return-supervised selection of a single initial configuration, and a
hybrid null-space controller that partitions execution between a learned
interior policy and a classical near-boundary regime. Beyond the system
itself, the experiments characterize the structure of the problem: the
candidate set, not the selection rule, bounds attainable performance
until enumeration replaces generation; complete closed-loop returns are
the only supervision under which static selection succeeds on a diverse
candidate set; execution-free selection has a measurable information
limit at roughly two thirds of the headroom; and the execution stage,
when trained with task-relative observations over randomized tasks,
requires no adaptation to the selection-induced distribution. The
system improves the average executed path substantially at an unchanged
budget of one selector inference and one rollout per task, and the
methodology---preregistered gates, sealed evaluation, and measured
rather than assumed couplings---is applicable beyond this task.

\bibliographystyle{IEEEtran}
\bibliography{references}

\end{document}
